\documentclass[11pt]{article}
\usepackage[margin=1in]{geometry}
\usepackage{amsmath,amssymb,amsthm}
\usepackage[hidelinks]{hyperref}
\usepackage{xurl}
\setlength{\emergencystretch}{2em}
\newtheorem{theorem}{Theorem}
\newtheorem{lemma}[theorem]{Lemma}
\newtheorem{corollary}[theorem]{Corollary}
\newcommand{\ar}{\operatorname{ar}}
\newcommand{\rk}{\operatorname{rank}}
\title{Rank-three augmentations for eight further alternating knots}
\author{Internal research working proof}
\date{17 November 2026}
\begin{document}\maketitle

Let
\[
\mathcal A_8=\{13a_{304},13a_{362},13a_{623},13a_{629},
13a_{3232},13a_{3259},13a_{3374},13a_{3375}\}.
\]
Each of these knots has an explicit rank-three symmetric cord matrix
and an irreducible three-dimensional KCH representation at meridian
parameter $-1$. The lower certificates below are independent of the
representation upper bounds and of the meridional-rank lower bounds.
Every statement about unrestricted augmentation rank uses the ordinary
$U=1$ cord algebra and the condition $\epsilon(\mu)\ne1$.

\begin{theorem}\label{thm:attainment}
For every $K\in\mathcal A_8$, there is an irreducible
three-dimensional representation of its knot group over a prime
cyclotomic field, taking meridians to reflections with eigenvalues
$-1,1,1$. The retained symmetric cord matrix has rank three.
Consequently $K$ admits a rank-three reflective augmentation and
$\ar(K,\mathbb C)\ge3$.
\end{theorem}

\section{The positive construction}

Write $B(u,v)$ for a symmetric form with $B(u,u)=2$ on every normal
vector. The associated reflection is
\[
                       R_u(v)=v-B(u,v)u.
\]
A symmetric cord matrix on a diagram with arcs $0,\ldots,12$ satisfies
\begin{equation}\label{eq:cord}
 c_{ii}=2,\qquad c_{\ell a}+c_{\ell b}=c_{\ell o}c_{oa}
 \quad\text{at every crossing }(o,a,b),\text{ for every }\ell.
\end{equation}
Here $o$ is the over-arc, and $a,b$ are the under-arcs.
The triples below are unoriented: their first entry is the over-arc,
and the last two entries may be interchanged.
Cornwell's diagrammatic cord presentation identifies these solutions
with evaluations of the reflective cord quotient
\cite[Theorem3.2 and Section3]{cornwell15}. Its rank is the matrix rank.
We also check the irreducible KCH representations directly, so the
unrestricted rank-three lower bound can instead use
\cite[Theorem1.2, Definition1.8 and Proposition2.12]{cornwell17}.

For the initial normal choices before the final sign adjustment, write
$A=u_0$, $C=u_j$ and $E=u_1$, where $j$ is given in the following
table. Let $\zeta$ be a primitive $p$th root of unity.
The numbers $b$ and the local type specify the symmetric form on the
ordered basis $(A,C,E)$:
\[
\begin{array}{c|c|c|c|c}
K&p&j&b=B(A,C)&\text{types in displayed order}\\\hline
13a_{304},13a_{3232}&37&6&\zeta^2+\zeta^{-2}&\mathrm I,\mathrm{II}\\
13a_{362},13a_{3259}&37&6&\zeta^{12}+\zeta^{-12}&\mathrm I,\mathrm{II}\\
13a_{623},13a_{3374}&7&8&\zeta^3+\zeta^{-3}&\mathrm I,\mathrm{II}\\
13a_{629},13a_{3375}&7&8&-(\zeta^3+\zeta^{-3})&\mathrm I,\mathrm{II}.
\end{array}
\]
Use
\begin{equation}\label{eq:gram}
 G_{\mathrm I}=\begin{pmatrix}2&b&1\\b&2&b\\1&b&2\end{pmatrix},
 \qquad
 G_{\mathrm{II}}=\begin{pmatrix}2&b&1\\b&2&0\\1&0&2\end{pmatrix}.
\end{equation}
Both determinants equal $2(3-b^2)$. They are nonzero in the specified
fields: this follows directly from the retained nonzero coefficient
vectors modulo $\Phi_p$, or because $b^2=3$ would force a nontrivial
$p$th root to have order dividing twelve.

Let $j'=j+1$. Before choosing the final signs of the normal vectors,
the five local crossings can be solved explicitly by
\[
\begin{array}{c|ccccc}
\text{type}&u_1&u_2&u_3&u_4&u_{j'}\\\hline
\mathrm I&E&E-A&E-A&-A&C\\
\mathrm{II}&E&E&E-A&-A&C.
\end{array}
\]
For type I, the local triples are
\[
 (0,1,2),(j,2,3),(3,0,1),(1,3,4),(2,j,j');
\]
for type II they are
\[
 (j,1,2),(0,2,3),(3,0,1),(1,3,4),(2,j,j').
\]
The other eight crossing triples are identical for the two knots
in each displayed pair; entries in the $K$ column are $13a$ subscripts:
\[
\begin{array}{c|llll}
K&\multicolumn{4}{c}{\text{remaining eight triples }(o,a,b)}\\\hline
304,3232&(4,8,9)&(11,4,5)&(5,11,12)&(9,5,6)\\
        &(7,9,10)&(12,7,8)&(8,10,11)&(10,12,0)\\[2pt]
362,3259&(4,10,11)&(8,4,5)&(5,8,9)&(11,5,6)\\
        &(7,9,10)&(12,7,8)&(9,11,12)&(10,12,0)\\[2pt]
623,3374&(4,6,7)&(10,4,5)&(5,11,12)&(9,5,6)\\
        &(6,12,0)&(7,10,11)&(11,7,8)&(12,9,10)\\[2pt]
629,3375&(4,10,11)&(11,4,5)&(5,7,8)&(9,5,6)\\
        &(6,12,0)&(10,6,7)&(7,11,12)&(12,9,10).
\end{array}
\]
Together with the appropriate five local triples, these give all
thirteen crossings used in \eqref{eq:cord}. The checker reconstructs
them from the literal PD and checks that the local and remaining
lists are disjoint and exhaustive.
Substitution in \eqref{eq:gram} verifies the five local reflection-line relations.
All these normals have norm two. In particular, the local lift adds
the new direction $E$ to the two-dimensional reduced core.

Here is a complete recipe for the other eight normals. Start with a
plane having norm-two basis $r_0,r_1$ and
$B(r_0,r_1)=q=\zeta+\zeta^{-1}$. Define
$r_{k+1}=qr_k-r_{k-1}$. The following table gives the eight initial
core roots for both knots in each pair. They span this plane; in
particular the displayed $A$ and $C$ are independent.
\[
\begin{array}{c|rrrrrrrr}
K&u_0&u_5&u_6&u_8&u_9&u_{10}&u_{11}&u_{12}\\\hline
304,3232&r_{17}&r_{22}&r_{15}&r_{34}&r_0&-r_{30}&-r_1&r_6\\
362,3259&r_{22}&r_{15}&r_{34}&r_0&-r_{30}&-r_1&r_6&r_{17}
\end{array}
\]
\[
\begin{array}{c|rrrrrrrr}
K&u_0&u_5&u_6&u_7&u_8&u_{10}&u_{11}&u_{12}\\\hline
623,3374&r_1&-r_1&-r_2&r_0&r_5&r_1&r_6&r_3\\
629,3375&-r_2&r_0&r_5&r_1&r_6&r_3&r_1&-r_1.
\end{array}
\]
All abbreviated knot subscripts here carry the prefix $13a$.
Extend the plane by $E$ using \eqref{eq:gram}, and insert the five
local roots from the preceding table. Thus every normal is specified
without solving an additional equation. These formulas were found
by lifting old reduced cores; their verification below uses every
crossing of the full new diagram.

Sign compatibility also requires a separate check. For each crossing,
the initially constructed roots satisfy $R_{u_o}u_a=\delta u_b$ with
$\delta\in\{1,-1\}$. The retained signs $s_i\in\{1,-1\}$ solve
$s_b=-\delta s_a$ on every under-arc pair. Replacing $u_i$ by
$s_iu_i$ therefore gives
\[
          u_a+u_b=B(u_o,u_a)u_o.
\]
All $169$ scalar equations in \eqref{eq:cord} are then verified
explicitly for each knot. This is stronger than a reflection-line
incidence check.
For reproducibility, the signs in arc order $0,1,\ldots,12$ are
\[
\begin{array}{c|rrrrrrrrrrrrr}
K&0&1&2&3&4&5&6&7&8&9&10&11&12\\\hline
304,3232&+&-&+&-&+&-&-&+&+&+&-&-&+\\
362,3259&+&-&+&-&+&-&-&+&+&-&-&+&+\\
623,3374&+&-&+&-&+&+&+&-&-&+&+&+&+\\
629,3375&+&-&+&-&+&+&+&+&+&-&-&-&+
\end{array}
\]
The input retains the resulting thirteen coordinate vectors, Gram
matrix, and full $13\times13$ cord matrix for each knot. The checker
does not import an old witness or a transfer theorem.

\section{Rank and irreducibility}

\begin{proof}[Proof of Theorem~\ref{thm:attainment}]
Let $V=\mathbb Q(\zeta)^3$ with the nondegenerate form
\eqref{eq:gram}, and take the signed retained roots $u_i$.
Their reflections have rank-one displacement and satisfy
all thirteen matrix Wirtinger relations. Both crossing signs have the
same involutory relation, so this defines a knot-group representation.
The principal cord minor on arcs $(0,j,1)$ has determinant
$2(3-b^2)\ne0$, unchanged by the final signs. Hence these three
roots and these three covectors $B(u_i,\cdot)$ are bases.
The full cord matrix factors through $V$ and has this nonzero minor,
so its rank is exactly three.

The pairing graph contains the full star centred at zero directly
from the printed formulas. Indeed, the recurrence gives
\[
                   B(r_a,r_k)=\zeta^{a-k}+\zeta^{-(a-k)}.
\]
This is nonzero: otherwise a power of the odd-order root $\zeta$
would have square $-1$ and hence order four. Since $A=\pm r_a$,
its pairing with every signed core root is nonzero. Its added local
pairings, before final signs, are $B(A,E)=1$,
$B(A,E-A)=-1$, and $B(A,-A)=-2$. Final signs preserve nonvanishing.
Thus the graph with an edge $i\!\sim\!j$ whenever
$B(u_i,u_j)\ne0$ contains all twelve edges $0\!\sim\!j$.
The checker also verifies this star exactly.
Let $W\subset V$ be a nonzero invariant subspace.
The common kernel of the covectors is zero, so some $w\in W$ and
some $i$ satisfy $B(u_i,w)\ne0$. Thus
\[
                    (I-R_{u_i})w=B(u_i,w)u_i
\]
puts $u_i$ in $W$. The same formula propagates roots along every
edge of the connected pairing graph. All thirteen roots therefore
belong to $W$, and they span $V$. This proves irreducibility.

The reflections are KCH meridian matrices at $\mu=-1$.
Cornwell's correspondence \cite{cornwell17} gives a rank-three ordinary
$U=1$ augmentation. Independently, \eqref{eq:cord} and the rank-three
minor give a reflective augmentation by the cord presentation
\cite{cornwell15}. The chosen embedding of the cyclotomic field into
$\mathbb C$ preserves every checked identity and nonzero minor.
\end{proof}

\begin{corollary}\label{cor:reflective}
For every $K\in\mathcal A_8$, its maximum complex reflective
augmentation rank is exactly three.
\end{corollary}
\begin{proof}
Theorem~\ref{thm:attainment} gives the lower bound. The upper bound is
the separate common-form reflection obstruction in
\path{UNIVERSAL_REFLECTION_OBSTRUCTION.tex}, Theorem
\texttt{thm:a6eight}, combined with the Gram-radical quotient construction
in \path{REFLECTIVE_AUGMENTATION_RANK_GAP.tex}, Lemma \texttt{lem:gram}.
Neither upper argument uses these positive matrices.
\end{proof}

\begin{corollary}\label{cor:augmentation}
For every $K\in\mathcal A_8$,
\[
                         \ar(K,\mathbb C)=3.
\]
\end{corollary}
\begin{proof}
The lower bound is Theorem~\ref{thm:attainment}. The separate generic
upper theorem in \path{GENERIC_KCH_EXTENSIONS.tex}, Theorem
\texttt{thm:all16} and Corollary \texttt{cor:augmentation}, bounds every
ordinary $U=1$ augmentation with $\epsilon(\mu)\ne1$ by rank three,
with arbitrary nonzero longitude parameter. That proof is independent
of this positive package.
\end{proof}

\section{Exact verification scope}

The input and standalone standard-library checker are
\path{a6_rank_three_augmentation_inputs.json} and
\path{verify_a6_rank_three_augmentations.py}. Each field is represented
as $\mathbb Q[X]/\Phi_p(X)$, where
$\Phi_p=1+X+\cdots+X^{p-1}$ for the prime $p=7$ or $37$.
The checker verifies primality and the Eisenstein coefficients of
$\Phi_p(X+1)$, which prove irreducibility over $\mathbb Q$.
It uses exact integer polynomial convolution followed by monic
$\Phi_p$ reduction; no sampling or floating-point roots are used.

For each knot the checker replays all thirteen $3\times3$ matrix
crossings, all $169$ cord equations, a nonzero principal minor,
root and covector spanning, and connectedness of the pairing graph.
It verifies the exact planar-map binding to a retained named census
diagram. The optional \texttt{--census} mode re-reads the pinned
independently packaged named PD and DT records. Corrupt cord matrices
and incorrect field polynomials are rejected. These checks are internal
exact computations, not formal verification or external specialist review.
The meridional-rank and bridge equalities, and every representation
upper bound, are separate dependencies.

\begin{thebibliography}{9}
\bibitem{cornwell17}
C. R. Cornwell, \emph{KCH representations, augmentations, and
$A$-polynomials}, Journal of Symplectic Geometry \textbf{15} (2017), 983--1017.
Theorem numbering is that of retained
\href{https://arxiv.org/abs/1310.7526v3}{arXiv:1310.7526v3}.
\bibitem{cornwell15}
C. R. Cornwell, \emph{Character varieties of knot complements and
branched double-covers via the cord ring}, retained
\href{https://arxiv.org/abs/1509.04962v1}{arXiv:1509.04962v1}.
\end{thebibliography}
\end{document}
